\section*{Chapter \ref{ch:b3:orogenic-wedges} --- Thrust Belts and Orogenic Wedges}

\begin{solution}{exo:b3:orogenic-wedges:1}
$\sin\ang{30} = 0.5$: $K = 2\times0.5/0.5 = 2$ and $\sigma'_1/\sigma'_3 = 1 +
K = 3$.
\end{solution}

\begin{solution}{exo:b3:orogenic-wedges:2}
$\phi = \arctan0.6 = \ang{31.0}$, $\sin\phi = 0.514$, $K = 2.12$. On a
horizontal base, $\alpha = 0.3/3.12 = \qty{0.0962}{rad} = \ang{5.5}$. On a
base dipping $\ang{3} = \qty{0.0524}{rad}$, $\alpha + \beta = 0.352/3.12 =
\qty{0.113}{rad} = \ang{6.5}$, so $\alpha = \ang{3.5}$.
\end{solution}

\begin{solution}{exo:b3:orogenic-wedges:3}
Depth $= (T - 10)/30$: apatite (U--Th)/He 1.3 to \qty{2.3}{km}; apatite fission
track 2.7 to \qty{3.7}{km}; zircon fission track $270/30 = \qty{9.0}{km}$. At
\qty{5}{km} the rock was at $10 + 150 = \qty{160}{\celsius}$: zircon
(U--Th)/He, which closes at 160 to \qty{200}{\celsius}, 5.0 to
\qty{6.3}{km} deep.
\end{solution}

\begin{solution}{exo:b3:orogenic-wedges:4}
Horizontal length $2/\tan\ang{30} = \qty{3.46}{km}$; along the dip
$2/\sin\ang{30} = \qty{4.0}{km}$. The bed is lifted by the ramp height,
\qty{2}{km}. The cut-off must travel the whole ramp: a slip of
\qty{4.0}{km}.
\end{solution}

\begin{solution}{exo:b3:orogenic-wedges:5}
$\phi = \arctan1.03 = \ang{45.8}$, $K = 5.08$; $\beta = \qty{0.105}{rad}$.
For $\lambda = 0.4$: $\alpha + \beta = (0.105 + 0.85\times0.6)/(1 +
5.08\times0.6) = 0.615/4.05 = \qty{0.152}{rad} = \ang{8.7}$, $\alpha =
\ang{2.7}$. For $\lambda = 0.9$: $0.190/1.51 = \qty{0.126}{rad} = \ang{7.2}$,
$\alpha = \ang{1.2}$. Raising the fluid pressure lowers the \omterm{def:b3:orogenic-wedges:critical-taper}{critical taper}:
the wedge becomes supercritical, slides forward on its base without
thickening, and its front advances until accretion and erosion bring it back
to the new taper.
\end{solution}

\begin{solution}{exo:b3:orogenic-wedges:6}
$S = A/r = 6.0/1.5 = \qty{4.0}{km}$. Read as a \omterm{def:b3:orogenic-wedges:thrust-folds}{detachment fold}, $S = A/h =
6.0/3.0 = \qty{2.0}{km}$: half the true value. The two readings must be told
apart by the shape of the fold and the presence of a ramp.
\end{solution}

\begin{solution}{exo:b3:orogenic-wedges:7}
$K = 2.12$ and $1 + K(1 - 0.4) = 2.27$. The equation gives $\mu_b(1 -
\lambda_b) = (\alpha + \beta)\times2.27 - \beta$, with $\beta =
\qty{0.0349}{rad}$. Jura: $0.0611\times2.27 - 0.0349 = 0.10$. South:
$0.218\times2.27 - 0.0349 = 0.46$. Byerlee at hydrostatic pressure:
$0.85\times0.6 = 0.51$. The base to the south is close to Byerlee's; the
Jura base is almost five times weaker, as expected of a décollement in
halite, whose friction is low, under fluids that are not overpressured.
\end{solution}

\begin{solution}{exo:b3:orogenic-wedges:8}
$\kappa = \qty{e-6}{m^2/s} = \qty{31.6}{km^2/Myr}$, and $v =
\qty{1}{km/Myr}$: $\mathrm{Pe} = 30/31.6 = 0.951$. The factor $0.951/(1 -
e^{-0.951}) = 1.55$ raises the gradient from 30 to \qty{46.5}{\celsius/km}.
For \qty{65}{\celsius}: $1 - e^{-\mathrm{Pe}\,z/L} = (55/900)(1 - e^{-0.951})
= 0.0375$, so $\mathrm{Pe}\,z/L = 0.0382$ and $z = \qty{1.21}{km}$, against
$55/30 = \qty{1.83}{km}$ without erosion.
\end{solution}

\begin{solution}{exo:b3:orogenic-wedges:9}
$\mathrm dh/\mathrm dt = \qty{0.8}{km}/\qty{1.2}{Myr} = \qty{0.67}{mm/yr}$.
With $f = 0.5$, $v = 0.67\times(1 - 0.5) = \qty{0.33}{mm/yr}$.
\end{solution}

\begin{solution}{exo:b3:orogenic-wedges:10}
With $\alpha = 0$ and $u = 1 - \lambda$: $\beta(1 + Ku) = \beta + \mu_b u$,
so $\beta Ku = \mu_b u$ and $\beta^* = \mu_b/K$, whatever $u$. For Taiwan,
$0.85/5.08 = \qty{0.167}{rad} = \ang{9.6}$. On a steeper base, the critical
surface slope is negative: any wedge with a level or foreland-sloping surface
is already supercritical, and it slides on its base without thickening. A
thrust belt cannot grow its normal shape there.
\end{solution}

\begin{solution}{exo:b3:orogenic-wedges:11}
As $\Lambda \to \infty$ and $v \to 0$, $m \to 0$, so $e^{-mz_c} \to 1$, and
the geotherm is straight, $G_0 = G(z_c) = G$: $f \to (G - \Gamma)/G = 1 -
\Gamma/G = 1 - 6.5/30 = 0.78$. The slope is then $v/(1 - f) = vG/\Gamma =
4.6v$. For $v = \qty{1}{mm/yr}$, $v/2\kappa = \qty{0.0158}{km^{-1}}$. At
\qty{5}{km}: $m = 0.0158 + \sqrt{0.0158^2 + 1.257^2} = \qty{1.273}{km^{-1}}$,
$e^{-1.273\times1.21} = 0.215$, $f = 40.0\times0.215/44.7 = 0.19$. At
\qty{20}{km}: $m = 0.0158 + 0.3146 = \qty{0.330}{km^{-1}}$, $e^{-0.399} =
0.671$, $f = 40.0\times0.671/44.7 = 0.60$.
\end{solution}

\begin{solution}{exo:b3:orogenic-wedges:12}
The section is a triangle of base $W$ and height $W\tan(\alpha + \beta)$:
$A = W^2\tan(\alpha + \beta)/2$. With $\mathrm dA/\mathrm dt = F_a$ and $A =
0$ at $t = 0$, $A = F_at$, which gives $W$. With $\tan\ang{6} = 0.105$:
$W = \sqrt{2\times40\times2/0.105} = \qty{39}{km}$ after \qty{2}{Myr}; $t =
60^2\times0.105/(2\times40) = \qty{4.7}{Myr}$ to reach \qty{60}{km}.
\end{solution}

\begin{solution}{pb:b3:orogenic-wedges:1}
\textbf{1.} $4.5 - 1.5 = \qty{3.0}{km}$ high; $3.0/\tan\ang{25} =
\qty{6.43}{km}$ long horizontally; $3.0/\sin\ang{25} = \qty{7.10}{km}$ along
the dip.

\textbf{2.} By the ramp height, \qty{3.0}{km}.

\textbf{3.} $S = A/r = 16.5/3.0 = \qty{5.5}{km}$.

\textbf{4.} The same, \qty{16.5}{km^2}: in a \omterm{def:b3:orogenic-wedges:thrust-folds}{fault-bend fold} the excess area
of every bed above the upper flat is $Sr$. As a \omterm{def:b3:orogenic-wedges:thrust-folds}{detachment fold}, $16.5/4.0 =
\qty{4.1}{km}$: a quarter too little.

\textbf{5.} $4\times(6.0 - 2.5) = \qty{14.0}{km}$, two and a half times the
\qty{5.5}{km} of slip on the frontal ramp.

\textbf{6.} The frontal horse, the one nearest the foreland: each new ramp
forms in the footwall ahead of the last. Each horse, carried up its ramp,
lifts the roof thrust and the cover above it as a \omterm{def:b3:orogenic-wedges:thrust-folds}{fault-bend fold} lifts its
hanging wall, so the roof is folded over the stack.

\textbf{7.} $\alpha = \arctan(2.1/60) = \ang{2.00}$; $\beta =
\arctan(4.2/60) = \ang{4.00}$; taper \ang{6.0}.

\textbf{8.} $2.4 + 8.2 = \qty{10.6}{km}$, less than about \qty{15}{km}: the
base is still frictional, and the Coulomb theory applies.

\textbf{9.} $\phi = \ang{45.8}$, $\sin\phi = 0.717$, $K = 5.08$. Dry:
$\alpha + \beta = (0.0699 + 0.85)/6.08 = \qty{0.151}{rad} = \ang{8.7}$,
$\alpha = \ang{4.7}$.

\textbf{10.} $\lambda = 1000/2500 = 0.40$: $(0.0699 + 0.51)/4.05 =
\qty{0.143}{rad} = \ang{8.2}$, $\alpha = \ang{4.2}$. The observed \ang{2.0}
would make the wedge strongly subcritical, thickening rather than advancing:
the fluids must be overpressured.

\textbf{11.} With $T = \qty{0.1049}{rad}$ and $u = 1 - \lambda$, $T(1 + Ku) =
\beta + \mu_bu$ gives $u = (\beta - T)/(TK - \mu_b) = (0.0699 -
0.1049)/(0.5327 - 0.85) = 0.110$: $\lambda = 0.89$.

\textbf{12.} $\sigma_v = 2500\times9.81\times10\,600 = \qty{260}{MPa}$; $P =
0.89\times260 = \qty{231}{MPa}$; hydrostatic $1000\times9.81\times10\,600 =
\qty{104}{MPa}$. The fluids carry \qty{127}{MPa} more than a column of
water would.

\textbf{13.} $(0.0699 + 0.85\times0.3)/(1 + 5.08\times0.3) = 0.325/2.52 =
\qty{0.129}{rad} = \ang{7.4}$, $\alpha = \ang{3.4}$. The wedge, at \ang{2.0},
is now subcritical: it stops advancing and thickens by out-of-sequence
thrusting behind its front until its surface slope reaches \ang{3.4}.

\textbf{14.} $F_a = \qty{10}{km/Myr}\times\qty{4.0}{km} = \qty{40}{km^2/Myr}$,
that is \qty{4.0e4}{m^3/yr} per kilometre of the belt.

\textbf{15.} $\bar e = 40/60 = \qty{0.67}{km/Myr} = \qty{0.67}{mm/yr}$.

\textbf{16.} $40\times100 = \qty{4000}{km^3/Myr} = \qty{4.0e6}{m^3/yr}$, or
$\qty{1.0e10}{kg/yr}$ (10 million tonnes a year).

\textbf{17.} $\mathrm dW/\mathrm dt = F_a/(W\tan(\alpha + \beta)) =
40/(60\times0.105) = \qty{6.3}{km/Myr}$.

\textbf{18.} No: at steady state the erosion must still remove the whole
incoming flux, so the mean rate stays \qty{0.67}{mm/yr}. But the rock added
beneath the rear is lifted there: erosion and exhumation would be fastest
over the rear, where the youngest cooling ages would be found.

\textbf{19.} $v = \qty{0.667}{km/Myr}$: $\mathrm{Pe} = 0.667\times30/31.6 =
0.634$; $G_0 = 30\times0.634/(1 - e^{-0.634}) = 30\times1.35 =
\qty{40.5}{\celsius/km}$.

\textbf{20.} $1 - e^{-\mathrm{Pe}\,z/L} = (55/900)(1 - e^{-0.634}) = 0.0287$,
so $\mathrm{Pe}\,z/L = 0.0291$ and $z_c = \qty{1.38}{km}$; $G(z_c) =
40.5\,e^{-0.0291} = \qty{39.3}{\celsius/km}$.

\textbf{21.} $1.38/0.667 = \qty{2.07}{Myr}$, the age of the sample at
\qty{2.00}{km}: at the mean elevation the deflection of the isotherm is zero.

\textbf{22.} $v/2\kappa = 0.667/63.1 = \qty{0.0106}{km^{-1}}$, $2\pi/\Lambda =
\qty{0.314}{km^{-1}}$: $m = 0.0106 + 0.3143 = \qty{0.325}{km^{-1}}$;
$e^{-0.325\times1.38} = 0.639$; $f = (40.5 - 6.5)\times0.639/39.3 = 0.552$.

\textbf{23.} The mean elevation is \qty{2.00}{km} and the mean age
\qty{2.068}{Myr}; $\sum\Delta t\,\Delta h = 0.823$ and $\sum\Delta t^2 =
0.552$, so the slope is \qty{1.49}{km/Myr}: an apparent rate of
\qty{1.49}{mm/yr}.

\textbf{24.} $v = 1.49\times(1 - 0.552) = \qty{0.667}{mm/yr}$.

\textbf{25.} The \omterm{def:b3:orogenic-wedges:thermochronology}{thermochronometer} and the flux balance agree: the wedge is
at steady state, and taking the apparent slope at face value would have
overestimated the rate more than twofold. The exhumation rate of the Kestrel
wedge is \textbf{\qty[text-series-to-math=true]{0.67}{mm/yr}}.
\end{solution}
